\documentclass[11pt]{article}
\usepackage[margin=1in]{geometry}
\usepackage{enumitem}
\usepackage{booktabs}
\usepackage{hyperref}
\usepackage{xcolor}
\usepackage{titlesec}
\setlength{\parindent}{0pt}
\setlength{\parskip}{0.5em}

\titleformat{\section}{\large\bfseries}{\thesection.}{0.5em}{}
\titleformat{\subsection}{\normalsize\bfseries}{\thesubsection}{0.5em}{}
\hypersetup{colorlinks=true, urlcolor=blue, linkcolor=black}

\title{\textbf{System Card: LFW Face Verification System}\\
\large MSML/MSAI 605 -- Milestone 4 (Final Release \texttt{v1.0-final})}
\author{}
\date{}

\begin{document}
\maketitle
\vspace{-2em}

\section{System Overview}
\textbf{What it does.} This system is a binary \emph{face-verification} model: given two
face images, it decides whether they belong to the \emph{same person} (label 1) or to
\emph{different people} (label 0), and emits a similarity score, a calibrated confidence
in $[0.5, 1.0]$, and a wall-clock latency.

\textbf{Pipeline.} The deployed system is a five-stage embedding-based pipeline:
\begin{enumerate}[noitemsep, leftmargin=*]
    \item \textbf{Preprocessing}: open RGB image, resize to $160\times 160$, normalize to $[-1, 1]$.
    \item \textbf{Embedding}: FaceNet (\texttt{InceptionResnetV1}, pretrained on VGGFace2) produces a 512-dim vector per image.
    \item \textbf{Similarity}: vectorized cosine similarity (\texttt{src/similarity.py}).
    \item \textbf{Decision}: $\text{score} \geq \theta \Rightarrow \text{same}$.
    \item \textbf{Confidence}: deterministic threshold-distance rule (\texttt{src/confidence.py}), $0.5$ at the boundary, $1.0$ at the score extremes.
\end{enumerate}

\textbf{Inputs / outputs.} Inputs: two image paths (or a CSV batch of pairs). Outputs:
\texttt{score}, \texttt{decision} ($\{$same, different$\}$), \texttt{confidence}, \texttt{latency\_ms}.

\textbf{Interfaces.} Importable Python (\texttt{src/inference.py::infer\_pair}); CLI
(\texttt{scripts/infer.py}); Dockerized CLI (\texttt{Dockerfile}, entry point
\texttt{python scripts/infer.py}).

\section{Intended Use}
\textbf{In scope.}
\begin{itemize}[noitemsep, leftmargin=*]
    \item Educational demonstration of an embedding-based verification pipeline.
    \item Offline batch verification on LFW-style \emph{cropped, frontal, well-lit} face crops.
    \item Coursework reproducibility, benchmarking, and ablations.
\end{itemize}

\textbf{Out of scope.}
\begin{itemize}[noitemsep, leftmargin=*]
    \item Any biometric, identification, or surveillance deployment.
    \item Liveness or anti-spoofing decisions; this system does no liveness check.
    \item Identity re-identification across very different sensor types or domains.
    \item Decisions about access, hiring, law enforcement, or any high-stakes context.
    \item Inputs that are not face crops (full scenes, multi-face images, non-face content).
\end{itemize}

\section{Data Summary}
\textbf{Source.} Labeled Faces in the Wild (LFW), \emph{deep-funneled} variant, accessed via
\texttt{kagglehub} (dataset \texttt{jessicali9530/lfw-dataset}, version 4).

\textbf{Splits.} Deterministic 70/15/15 train/val/test split with \texttt{seed=42}; identities
stay within a single split. Pair generation also uses \texttt{seed=42}. After filtering
unreadable image paths the released system uses 4{,}199 train, 900 val, and 901 test pairs
(see \texttt{outputs/pairs/}). \emph{Validation} is the only split used for threshold
selection; \emph{test} is held out.

\textbf{Important data limitations.}
\begin{itemize}[noitemsep, leftmargin=*]
    \item LFW is web-scraped from news photography. It is heavily skewed toward
    public figures (politicians, athletes, celebrities) and toward adult, frontal,
    well-lit faces.
    \item Demographic composition of LFW is uneven and not labeled in this dataset, so
    the released system carries no reliable per-group metric.
    \item Some identities dominate the pair set (e.g.\ George\_W\_Bush). Milestone~2
    introduced an optional identity-cap + label-rebalance pass
    (\texttt{outputs/pairs\_improved/}); the released system reports metrics on the
    \emph{original} pair set so the numbers are comparable across milestones.
\end{itemize}

\textbf{Path portability.} The generated pair / score CSVs in \texttt{outputs/pairs/}
and \texttt{outputs/scores/} contain absolute paths into the local kagglehub cache
and are therefore machine-specific. They are intentionally gitignored. The
repository ships the deterministic generators, \texttt{configs/m1.yaml}, and
\texttt{seed=42}; the reproducibility checklist regenerates the CSVs against the
grader's machine. Because the seed and sort order are fixed, the resulting pair
content is identical across machines (only the path prefix differs).

\section{Operating Threshold and Key Metrics}
\textbf{Threshold selection.} Sweep cosine-similarity thresholds on the validation split
over $[-0.5, 1.0]$ at 200 evenly spaced points; pick the threshold that maximizes
\textbf{F1} on validation. The threshold is locked before touching test.

\textbf{Selected operating threshold:} $\boldsymbol{\theta = 0.351759}$
(see \texttt{outputs/threshold/selected\_threshold.json}, also written into
\texttt{configs/m3.yaml::inference.threshold}).

\begin{table}[h]
\centering
\small
\begin{tabular}{lrrrr}
\toprule
\textbf{Split} & \textbf{Accuracy} & \textbf{Precision} & \textbf{Recall} & \textbf{F1} \\
\midrule
Validation ($N=900$, 457 pos) & 0.9633 & 0.9670 & 0.9606 & 0.9638 \\
Test ($N=901$, 468 pos)       & 0.9734 & 0.9723 & 0.9765 & 0.9744 \\
\bottomrule
\end{tabular}
\end{table}

\textbf{Test confusion matrix} ($\theta=0.351759$): TP=457, FP=13, TN=420, FN=11
(\texttt{outputs/eval/test\_confusion\_matrix.json}).

\textbf{Confidence.} Calibrated to $[0.5, 1.0]$ from
$\frac{|score - \theta|}{\text{range to extreme}}$, so $0.5$ means ``borderline'' and
$1.0$ means ``score is at the extreme of its range.'' This is a deterministic mapping
with no learned parameters; it does not produce calibrated probabilities of identity.

\section{Failure Modes and Limitations}
The verifier was observed to be less reliable, or can break, on the following inputs:

\begin{itemize}[noitemsep, leftmargin=*]
    \item \textbf{Visually similar different identities} -- e.g.\ relatives, makeup, or
    near-look-alikes. Cosine similarity is high; the decision can flip to ``same.''
    \item \textbf{Off-pose / non-frontal faces} -- LFW is largely frontal; profile and
    extreme yaw inputs are out-of-distribution and produce noisier scores.
    \item \textbf{Strong lighting or low-quality crops} -- under/overexposed images,
    motion blur, low resolution below the 160$\times$160 expected size.
    \item \textbf{Partial occlusion} -- masks, sunglasses, hands, hair covering the face.
    \item \textbf{Different sensor / domain} -- non-RGB, infrared, heavily filtered, or
    cartoon images are not in the training distribution and outputs should not be trusted.
    \item \textbf{Non-face inputs} -- the pipeline does not detect or reject
    non-face content; it will silently produce a similarity score.
    \item \textbf{Dataset-specific bias} -- LFW is overrepresented by adult public
    figures; performance on under-represented populations is not characterized here
    and should be assumed worse than the headline numbers.
\end{itemize}

\section{Fairness-Risk Discussion}
This system uses a public dataset (LFW) without reliable demographic metadata, so we do
\emph{not} make subgroup performance claims. We instead document risk categories that
peer-reviewed face-recognition literature (e.g.\ NIST FRVT, Buolamwini \& Gebru 2018) has
repeatedly shown to drive uneven accuracy:

\begin{itemize}[noitemsep, leftmargin=*]
    \item \textbf{Population coverage}: LFW skews white, male, and adult. A
    pre-trained VGGFace2 backbone was trained on a similarly biased web corpus.
    Performance on women, on darker-skinned individuals, and on children should be
    expected to be \emph{worse} than the aggregate metrics in this card.
    \item \textbf{Image-quality differences}: image quality often correlates with
    demographics (e.g.\ camera availability, lighting in source photos). Score
    distributions may therefore diverge along uneven lines even when the model is
    ``the same.''
    \item \textbf{Threshold selection}: the operating threshold is tuned to maximize
    aggregate F1 on validation. A single global threshold can disadvantage subgroups
    where the score distribution shifts.
    \item \textbf{Deployment misuse}: the system is \emph{not} a re-identification
    or surveillance tool. Using it as one would amplify any subgroup-level error
    against the people most affected.
\end{itemize}

\textbf{Mitigations in this release.}
(i) The Out-of-scope section above explicitly forbids high-stakes use.
(ii) The system emits a calibrated confidence so callers can apply abstention near
$\theta$.
(iii) The data-centric improvement script (\texttt{scripts/apply\_data\_centric.py})
demonstrates how to cap dominant identities for fairer evaluation; the metric numbers
in this card are reported on the unfiltered pair set so they remain comparable to
Milestone 2.

\section{Operational Constraints}
\begin{itemize}[noitemsep, leftmargin=*]
    \item \textbf{Input format}: RGB face crops; the system internally resizes to
    $160\times 160$ and normalizes to $[-1, 1]$.
    \item \textbf{Hardware}: tested on CPU only (Apple~M-series, 8 logical cores).
    GPU is supported in code but optional; see profiling report for measurement
    methodology.
    \item \textbf{Latency / throughput} (CPU baseline, see profiling report):
    end-to-end median $\approx 124$\,ms per pair at batch~1; throughput peaks near
    $\approx 12$ pairs/s at batch~8 and degrades sharply at batch~$\geq 16$ on the
    8-core profiling host (likely thermal / memory pressure).
    \item \textbf{Embedding stage dominates runtime} ($\approx 94\%$ of single-pair
    end-to-end). Preprocessing is $\approx 6\%$; cosine scoring is $<0.2\%$.
    \item \textbf{Determinism}: all random operations use \texttt{seed=42}. The
    inference path itself has no randomness.
    \item \textbf{Dependencies}: pinned in \texttt{requirements.txt}; the
    \texttt{Dockerfile} bakes the FaceNet weights into the image so the container
    works offline.
    \item \textbf{Concurrency}: 4-worker thread load test from Milestone~3 hit
    $\approx 22.6$ req/s with stable P95 (\texttt{outputs/load\_test/runtime\_summary.json}).
\end{itemize}

\section{Reproducibility Pointer}
\begin{itemize}[noitemsep, leftmargin=*]
    \item \textbf{Repository}: this directory; final tag \texttt{v1.0-final}.
    \item \textbf{Reproducibility checklist}: \texttt{reports/reproducibility\_checklist.md}.
    \item \textbf{Profiling report}: \texttt{reports/profiling\_report.pdf} (sources in \texttt{reports/profiling\_report.tex}; raw artifacts in \texttt{outputs/profiling/}).
    \item \textbf{README}: top-level \texttt{README.md} -- copy-pastable setup, CLI, Docker, and profiling commands.
    \item \textbf{Final config}: \texttt{configs/m3.yaml} (seed, threshold $0.351759$, score range, preprocessing).
    \item \textbf{Selected threshold artifact}: \texttt{outputs/threshold/selected\_threshold.json}.
    \item \textbf{Test metrics artifact}: \texttt{outputs/eval/test\_metrics.json}.
\end{itemize}

\end{document}
